
\color{uniblue}\subsection[Выбор уставок ДЗТ]{Выбор уставок ДЗТ} 
\color{black}

\begin{enumerate}

\item
Параметры из раздела <<Общие уставки>> не требуют расчетов, принимаются следующие значения с учетом исходных данных и заданных параметров:\\
\textit{Сторона 1: 1} = Предусмотрено;\\
\textit{Сторона 2: 0} = Не предусмотрено;\\
\textit{Сторона 3: 1} = Предусмотрено;\\
\textit{Схема соединения трансформаторов тока стороны 1: 0} = «Звезда»;\\
\textit{Схема соединения трансформаторов тока стороны 3: 0} = «Звезда»;\\
\textit{Базисная мощность: 16,0} МВА.
\item
Дополнительные параметры приведения сторон:\\
\textit{Номинальное напряжение стороны 1: 37,0} кВ;\\
\textit{Номинальный первичный ток трансформатора тока стороны 1: 400,0} А;\\
\textit{Номинальный вторичный ток трансформатора тока стороны 1: 5,0} А;\\
\textit{Компенсация токов 3I0 для стороны 1: 0} = Без компенсации;\\
\textit{Векторная группа стороны 1: 0};\\
\textit{Номинальное напряжение стороны 3: 6,3} кВ;\\
\textit{Номинальный первичный ток трансформатора тока стороны 3: 1500,0} А;\\
\textit{Номинальный вторичный ток трансформатора тока стороны 3: 5,0} А;\\
\textit{Компенсация токов 3I0 для стороны 3: 0} = Без компенсации;\\
\textit{Векторная группа стороны 3: 11}.

\item Промежуточный коэффициент амплитудного выравнивания для стороны ВН :
\begin{equation*}
    k_{\textsl{ампл, ВН}} = \frac{I_{\textsl{перв,k}}}{  I_{\textsl{втор,k}} \cdot  k_{\textsl{сх}}   \cdot I_{\textsl{баз,k}}}  = \frac{400}{5\cdot 1 \cdot 249,68} = 0,3204.
\end{equation*}

\needspace{3\baselineskip}
\item Промежуточный коэффициент амплитудного выравнивания для стороны НН:
\begin{equation*}
    k_{\textsl{ампл, НН}} = \frac{I_{\textsl{перв,k}}}{  I_{\textsl{втор,k}} \cdot  k_{\textsl{сх}}   \cdot I_{\textsl{баз,k}}}  = \frac{1500}{5\cdot 1 \cdot 1466,3} = 0,2046.
\end{equation*}

\item 
Характеристика срабатывания формируется из трех зон: начального горизонтального участка и двух наклонных участков торможения.

\item \textbf{Расчет начального дифференциального тока срабатывания}

Дифференциальный ток срабатывания начального участка отстраивается от режима нормальной номинальной нагрузки. Ток небаланса вычисляется по выражению (\ref{eq:dzt1}):

\begin{equation*}
    I_{\textsl{нб.расч.}} =  I_{\textsl{нб.расч.*}}^{'} + I_{\textsl{нб.расч.*}}^{''} + I_{\textsl{нб.расч.*}}^{'''} = 0,1 + 0,12 + 0,03 = 0,25\; (\textsl{о.е.}),
\end{equation*}
где:
\begin{itemize}
    \item составляющая, обусловленная погрешностью ТТ (\ref{eq:dzt2}):
{\begin{equation*}
    I_{\textsl{нб.расч.*}}^{'}=  k_{\textsl{пер}} \cdot k_{\textsl{одн}} \cdot \epsilon \cdot I_{\textsl{расч}} = 1,0 \cdot 1,0 \cdot 0,1 \cdot 1,0 = 0,1\; (\textsl{о.е.}),
\end{equation*}}
    \item составляющая небаланса, вносимая устройством регулирования напряжения трансформатора (\ref{eq:dzt3}):
\begin{equation*}
    I_{\textsl{нб.расч.*}}^{''}= \Delta U_{\textsl{рег}} \cdot I_{\textsl{расч}} = 0,12 \cdot 1,0 = 0,12\; (\textsl{о.е.}),
\end{equation*}
\begin{equation*}
\text{где\quad}  
\Delta U_{\textsl{рег}} = \frac {|8 \cdot 1,5\%|+|8 \cdot (-1,5\%)|}{2 \cdot 100 \%} =0,12\; (\textsl{о.е.}),
\end{equation*}
\item составляющая, обусловленная погрешностью выравнивания токов плеч (\ref{eq:dzt4}):
\begin{equation*}
    I_{\textsl{нб.расч.*}}^{'''}= f_{\textsl{выр}} \cdot I_{\textsl{расч}} = 0,03 \cdot 1,0 = 0,03\; (\textsl{о.е.}),
\end{equation*}
\end{itemize}
По выражению (\ref{eq:dzt5}) определяем параметр <<Дифференциальный ток срабатывания начального участка>>:
\begin{equation*}
    I_{\textsl{ср}}= k_{\textsl{отс}} \cdot I_{\textsl{нб.расч}}=1,2 \cdot 0,25 = 0,3\; (\textsl{о.е.}) ,
\end{equation*}
Параметр \textbf{$I_{\textsl{ср(ДТЗт)}}$} <<Дифференциальный ток срабатывания начального участка>> принимается \textbf{равным~0,3}.


\item \textbf{Расчет коэффициента торможения первого наклонного участка}

\item 
Определение тока небаланса по выражению (\ref{eq:dzt1}) для выбора параметра <<Коэффициент торможения первого наклонного участка>>. Расчетной точкой является режим максимального внешнего КЗ на стороне НН ($I^{(3)}_{\textsl{K2,max, ВН}} = 1592,3~\textsl{А}$). Относительный ток расчетного режима:

\begin{equation*}
    I_{\textsl{расч.}} =  \frac {1592,3\cdot 1 \cdot 0,3204 \cdot5}{400} = 6,38\; (\textsl{о.е.}),
\end{equation*}

Определение тока небаланса:

\begin{equation*}
    I_{\textsl{нб.расч.}} =  I_{\textsl{нб.расч.*}}^{'} + I_{\textsl{нб.расч.*}}^{''} + I_{\textsl{нб.расч.*}}^{'''} = 2,164 + 0,765 + 0,324 = 2,232\; (\textsl{о.е.}),
\end{equation*}
где:

\begin{itemize}

\item составляющая, обусловленная погрешностью ТТ (\ref{eq:dzt2}):
{\begin{equation*}
    I_{\textsl{нб.расч.*}}^{'}=  k_{\textsl{пер}} \cdot k_{\textsl{одн}} \cdot \epsilon \cdot I_{\textsl{расч}} = 2 \cdot 1,0 \cdot 0,1 \cdot 6,38 = 1,276\; (\textsl{о.е.}),
\end{equation*}}

\item составляющая небаланса, вносимая устройством регулирования напряжения трансформатора (\ref{eq:dzt3}):
\begin{equation*}
    I_{\textsl{нб.расч.*}}^{''}= \Delta U_{\textsl{рег}} \cdot I_{\textsl{расч}} = 0,12 \cdot 6,38 = 0,765\; (\textsl{о.е.}),
\end{equation*}

\item составляющая, обусловленная погрешностью выравнивания токов плеч (\ref{eq:dzt4}):
\begin{equation*}
    I_{\textsl{нб.расч.*}}^{'''}= f_{\textsl{выр}} \cdot I_{\textsl{расч}} = 0,03 \cdot 6,38 = 0,191\; (\textsl{о.е.}),
\end{equation*}

\end{itemize}
По выражению (\ref{eq:dzt5}) определяем параметр <<Дифференциальный ток срабатывания начала второго участка>>:
\begin{equation*}
    I_{\textsl{ср(т1)}}= k_{\textsl{отс}} \cdot I_{\textsl{нб.расч}}=1,2 \cdot 2,232 = 2,678\; (\textsl{о.е.}) ,
\end{equation*}

\needspace{3\baselineskip}
По выражению (\ref{eq:dzt6}) определяем параметр <<Коэффициент торможения первого наклонного участка>>:
\begin{equation*}
    k_{\textsl{т1}} = \frac {I_{\textsl{ср(т2)}} - I_{\textsl{ср}}}{I_{\textsl{т2}}-I_{\textsl{т1}}} = \frac {2,678 - 0,3}{6,38-1,0}=0,44\; (\textsl{о.е.}),
\end{equation*}
\item 
Параметр \textbf{$k_{\textsl{т1}}$} <<Коэффициент торможения первого наклонного участка>> принимается \textbf{равным~0,44}.

\item 
Параметр \textbf{$I_{\textsl{т1}}$} <<Ток торможения, соответствующий началу первого наклонного участка>> принимается \textbf{равным~1,0}. 


\item \textbf{Расчет коэффициента торможения второго наклонного участка}


\item 
Определение тока небаланса по выражению (\ref{eq:dzt1}) для выбора параметра <<Коэффициент торможения второго наклонного участка>>.

Для расчета $k_{\textsl{т2}}$  используем минимальную предельную кратность, при которой ТТ войдет в насыщение (\ref{eq:dzt71}).

Приведенная кратность ТТ ВН:
\begin{equation}
    k'_{\textsl{10,ВН}} = 30 \cdot \frac{400}{249,68}=48,06\; (\textsl{о.е.}).
\end{equation}

Приведенная кратность ТТ НН:
\begin{equation}
    k'_{\textsl{10,НН}} = 15 \cdot \frac{1500}{1466,3}=15,34\; (\textsl{о.е.}).
\end{equation}

\item 
Расчетной принимается наименьшая кратность: $I_{\textsl{расч.}}=15,34$ о.е.

\item Определение тока небаланса

\begin{equation*}
    I_{\textsl{нб.расч.}} =  I_{\textsl{нб.расч.*}}^{'} + I_{\textsl{нб.расч.*}}^{''} + I_{\textsl{нб.расч.*}}^{'''} = 6,136 + 1,841 + 0,46 = 8,437\; (\textsl{о.е.}),
\end{equation*}
где:
\begin{itemize}
    \item составляющая, обусловленная погрешностью ТТ (\ref{eq:dzt2}):
{\begin{equation*}
    I_{\textsl{нб.расч.*}}^{'}=  k_{\textsl{пер}} \cdot k_{\textsl{одн}} \cdot \epsilon \cdot I_{\textsl{расч}} = 4,0 \cdot 1,0 \cdot 0,1 \cdot 15,34 = 6,136\; (\textsl{о.е.}),
\end{equation*}}

    \item составляющая небаланса, вносимая устройством регулирования напряжения трансформатора (\ref{eq:dzt3}):
\begin{equation*}
    I_{\textsl{нб.расч.*}}^{''}= \Delta U_{\textsl{рег}} \cdot I_{\textsl{расч}} = 0,12 \cdot 15,34 = 1,841\; (\textsl{о.е.}),
\end{equation*}

\item составляющая, обусловленная погрешностью выравнивания токов плеч (\ref{eq:dzt4}):
\begin{equation*}
    I_{\textsl{нб.расч.*}}^{'''}= f_{\textsl{выр}} \cdot I_{\textsl{расч}} = 0,03 \cdot 15,34 = 0,46\; (\textsl{о.е.}),
\end{equation*}

\end{itemize}

\item
Значение расчетного дифференциального тока, соответствующего максимальной предельной кратности ТТ:
\begin{equation*}
    I_{\textsl{ср.max}}= k_{\textsl{отс}} \cdot I_{\textsl{нб.расч}} = 1,2 \cdot 8,437 = 10,124\; (\textsl{о.е.}),
\end{equation*}

\item
Определяем относительный расчетный дифференциальный ток, соответствующий концу первого участка тормозной характеристики:
\begin{equation*}
    I_{\textsl{ср.(т2)}}= I_{\textsl{ср.(т1)}} +   k_{\textsl{т1}} \cdot (I_{\textsl{т1}}-I_{\textsl{т2}}) = 0,3 +0,44 \cdot(3-1) = 1,18\; (\textsl{о.е.}),
\end{equation*}

\item 
По выражению (\ref{eq:dzt7}) определяем параметр <<Коэффициент торможения второго наклонного участка>>:
\begin{equation*}
    k_{\textsl{т2}} = \frac {I_{\textsl{ср.макс}} - I_{\textsl{ср(т2)}}}{I_{\textsl{т2макс}}-I_{\textsl{т2}}} = \frac {10,124 - 1,18}{15,34-3,0}=0,72\; (\textsl{о.е.}),
\end{equation*}

\item 
Параметр \textbf{$k_{\textsl{т2}}$} <<Коэффициент торможения второго наклонного участка>> принимается \textbf{равным~0,72}.

\item 
Параметр \textbf{$I_{\textsl{т2}}$} <<Ток торможения, соответствующий началу второго наклонного участка>> принимается \textbf{равным~3,0}.



\item Проверка чувствительности защиты

Расчетным режимом является минимальное внутреннее повреждение: двухфазное металлическое КЗ на выводах 6,3 кВ ($I^{(2)}_{\textsl{K2,max, НН}} = 769,9~\textsl{А}$).

\item
Относительный ток расчетного режима:
\begin{equation*}
    I_{\textsl{расч}} = \frac {769,9 \cdot 1 \cdot 0,3204 \cdot 5}{400} = 3,083\; (\textsl{о.е.}).
\end{equation*}

\item
Коэффициент чувствительности определяется по выражению (\ref{eq:dzt92})
\begin{equation*}
    k_{\textsl{ч}} = \frac {3,083 - 0,43 \cdot(0,5 \cdot 3,083 - 1,0)}{0,29}=9,83,
\end{equation*}

Вычисляем ординату пересечения (\ref{eq:dzt921}):
\begin{equation*}
    I_{\textsl{дифф}} = \frac {3,083}{9,83} = 0,314\; (\textsl{о.е.}).
\end{equation*}

Сравниваем с границей участка (\ref{eq:dzt922}):
\begin{equation*}
    0,314 \leq 0,29 + 0,43 \cdot (3,0 - 1,0),
\end{equation*}
\begin{equation*}
    0,314 \leq 1,15.
\end{equation*}

\item 
Условие истинно. Расчет корректен.

\item 
Дифференциальная токовая отсечка

Определение тока небаланса по выражению для выбора параметра «Дифференциальный ток срабатывания отсечки» исходит из условия отстройки от тока небаланса при максимальном сквозном коротком замыкании.

Расчетной точкой является режим максимального внешнего КЗ на стороне НН ($I^{(3)}_{\textsl{K2,max, ВН}} = 1592,3~\textsl{А}$). 
\item
Относительный ток расчетного режима:
\begin{equation*}
    I_{\textsl{расч}} = \frac {1592,3 \cdot 1 \cdot 0,3204 \cdot 5}{400} = 6,38\; (\textsl{о.е.}).
\end{equation*}

Определение тока небаланса по выражению (\ref{eq:dzt1}) для выбора параметра <<Дифференциальный ток срабатывания отсечки>>:

\begin{equation*}
    I_{\textsl{нб.расч.}} =  I_{\textsl{нб.расч.*}}^{'} + I_{\textsl{нб.расч.*}}^{''} + I_{\textsl{нб.расч.*}}^{'''} = 2,552 + 0,765 + 0,191 = 3,5\; (\textsl{о.е.}),
\end{equation*}
где:
\begin{itemize}
    \item составляющая, обусловленная погрешностью ТТ (\ref{eq:dzt2}):
{\begin{equation*}
    I_{\textsl{нб.расч.*}}^{'}=  k_{\textsl{пер}} \cdot k_{\textsl{одн}} \cdot \epsilon \cdot I_{\textsl{расч}} = 4,0 \cdot 1,0 \cdot 0,1 \cdot 6,38 = 2,552\; (\textsl{о.е.}),
\end{equation*}}
    \item составляющая небаланса, вносимая устройством регулирования напряжения трансформатора (\ref{eq:dzt3}):
\begin{equation*}
    I_{\textsl{нб.расч.*}}^{''}= \Delta U_{\textsl{рег}} \cdot I_{\textsl{расч}} = 0,12 \cdot 6,38 = 0,765\; (\textsl{о.е.}),
\end{equation*}
\item составляющая, обусловленная погрешностью выравнивания токов плеч (\ref{eq:dzt4}):
\begin{equation*}
    I_{\textsl{нб.расч.*}}^{'''}= f_{\textsl{выр}} \cdot I_{\textsl{расч}} = 0,03 \cdot 6,38 = 0,191\; (\textsl{о.е.}),
\end{equation*}
\end{itemize}
По выражению (\ref{eq:dzt5}) определяем параметр <<Дифференциальный ток срабатывания отсечки>>:
\begin{equation*}
    I_{\textsl{ср(ДТО)}}= k_{\textsl{отс}} \cdot I_{\textsl{нб.расч}}=1,5 \cdot 3,5 = 5,262\; (\textsl{о.е.}) ,
\end{equation*}

Поскольку значение параметра <<Дифференциальный ток срабатывания отсечки>> по условиям отстройки от броска тока намагничивания по выражению (\ref{eq:dzt16}) ($I_{\textsl{ср(ДТО)}} \geq 6,0$) оказалось выше, чем по условию отстройки от тока небаланса при внешнем КЗ в максимальном режиме ($I_{\textsl{ср(ДТО)}}=5,262$), принимаем параметр \textbf{$I_{\textsl{ср(ДТО)}}$} <<Дифференциальный ток срабатывания отсечки>> \textbf{равным 6,0}.

\end{enumerate}